\section{Introduction}

Why have stock-price reactions to index changes weakened as index investing has grown? \citet{greenwood2025} examine this question for the S\&P 500 and document a substantial decline in the returns associated with additions and deletions. Their published abstract reports average addition returns below 1\% and deletion returns of 0.1\% in their latest decade. This finding matters because an index change alters the portfolio that benchmark-tracking investors seek to hold. A smaller price reaction can therefore reflect changes in how markets accommodate predictable trading, rather than a smaller role for index investing itself. Distinguishing these possibilities requires examining trading alongside returns.

This paper replicates and extends the empirical approach of Greenwood and Sammon in the Swedish equity market. Our research question is: \emph{How has the return pattern around OMX Stockholm Benchmark reconstitutions changed, and how does that pattern relate to measured benchmark-linked trading?} We first examine the evolution of announcement, implementation and total-window returns. We then consider announcement timing, the characteristics of affected securities, reported fund holdings and the concentration of trading around implementation. This sequence follows the economic question of the original study while adapting the measures to the Swedish index and available data.

The Swedish setting is informative for two reasons. First, OMX Stockholm Benchmark, hereafter OMXSB, selects constituents with reference to trading activity and industry representation; membership does not follow a single market-capitalization cutoff \citep{nasdaq2013,nasdaq2026}. Its scheduled reviews allow us to distinguish the announcement of a change from its implementation. Second, fund-level holdings allow us to examine the quantity of measured benchmark-linked trading alongside the return response. We supplement the commercial holdings data with reports from Finansinspektionen, the Swedish Financial Supervisory Authority. These features permit an extension beyond a comparison of national average returns: we can examine whether price and trading patterns tell a consistent story within the Swedish setting.

Our event inventory covers 39 reviews from July 2007 through June 2026. The main return sample contains 256 additions and 215 deletions. For each event, we calculate market-adjusted returns relative to the OMXSB gross total-return index. The main window begins one trading day before the announcement and ends one trading day after the change becomes effective. We report four historical blocks, together with comparisons between them and pooled estimates for the later decade. Uncertainty estimates account for the fact that several securities enter or leave at the same review. The holdings analysis uses its own matched sample because historical holdings are not available for every return observation.

Three findings emerge. First, the Swedish return pattern shows lower recent averages than in the earliest years, but no smooth disappearance. Average addition returns are 8.52\%, 0.25\%, 2.85\% and 1.50\% in 2007--2011, 2012--2016, 2017--2021 and 2022--2026, respectively. Pooling 2017--2026 gives an average of 2.15\%, with a 95\% confidence interval of [1.01\%, 3.39\%]. For 2022--2026 alone, the interval is [-0.10\%, 3.22\%]. The latest observations therefore leave substantial uncertainty about the remaining return, rather than establishing that it is economically negligible. Deletion returns are less precisely estimated: the later-decade mean is -1.04\%, with an interval spanning zero.

Second, the time between announcement and implementation has changed. Across the reviews contributing additions, the median gap falls from fourteen trading days in 2007--2011 to six in 2022--2026. The full event window consequently covers different durations. In an exploratory timing extension, we compare a common seven-trading-day window, from the day before announcement through five trading days after it. This window ends before implementation for every event, and each period retains exactly the same stocks as its main return estimate. Average addition returns over these seven days are 1.19\% in the earliest period and 1.80\% in the latest, compared with 8.52\% and 1.50\% over the full windows. The large early-to-late difference is therefore concentrated after the common initial window. This finding makes the time available for returns to accumulate an important part of the historical comparison.

Third, the return response coincides with measurable benchmark-linked trading. Among 136 additions with matched holdings data in 2017--2026, the mean event-window return is 2.71\% and measured net purchases represent 1.03\% of shares outstanding. The ratio of those means is 2.63 percentage points of return per percentage point of measured ownership change. The Finansinspektionen comparison provides a check on the scale of the holdings measure: in a matched sample of 106 later additions, the reported ownership changes are 1.14\% and 1.05\% under the two reporting sources. Historical comparisons are less secure because observed investor coverage changes and the earliest holdings samples are small.

Trading also becomes more concentrated immediately before implementation. For additions, the final trading day before the effective date accounts for 9.44\% of the surrounding 45-day trading volume in the earliest period and 15.24\% in the latest. Deletions display a similar increase. Relatively small recent implementation-window returns therefore coexist with concentrated trading around the date on which benchmark portfolios change. The amount and timing of trading need not move proportionally with the contemporaneous return. Daily volume identifies this concentration of activity without distinguishing closing-auction trades from other transactions.

The timing extension also follows a fixed event sample through twenty trading days after implementation. For additions in 2017--2026, the cumulative market-adjusted return over this longer window is 3.30\%, with a 95\% interval of [1.31\%, 5.28\%]. The positive result remains under separate exclusions of an influential event and simultaneous MSCI changes. The latest-period estimate is more sensitive to those exclusions. Together, the common initial window and the longer horizon show why the timing of the measured return matters for interpreting the headline averages.

Our contribution is a Swedish replication that connects the return comparison to announcement timing, reported holdings and implementation trading. A positive addition return remains visible over the later decade, while the shorter recent announcement-to-implementation gap changes the interpretation of the historical comparison. These findings are observational: they do not establish that the calendar caused the change in returns, that the initial-window means are equivalent, or that index demand causes a permanent price change. Selection, firm news and changing holdings coverage remain possible influences. The return-to-holdings ratio is consequently descriptive rather than an identified demand elasticity. The economic contribution is evidence on which return and trading patterns extend to Sweden and how their measurement affects the answer.

Section 2 positions the study in the related literature. Section 3 describes the index, data and empirical methods. Section 4 presents the replication results, followed by their interpretation and the extension in Section 5. Section 6 concludes.
